\section{Resultados e discussão}\label{sec:resultados}

\subsection{O que mudou em relação ao trabalho anterior}

O trabalho anterior\footnote{VICTOR, C. \textit{Prova P1 de Finanças do Setor Público} (matrícula 231020069). 2026. Repositório GitHub: \url{https://github.com/cauaviictor/P1---FSP---231020069---CAUA-VICTOR}, \textit{commit} 5209bcec.} disponibilizou base editável (planilha com fórmulas e arquivos do Finbra), preservada integralmente na pasta \texttt{01\_base\_original} do repositório. A conferência dessa base orientou as correções e ampliações resumidas no Quadro~\ref{tab:mudancas}; o registro detalhado está em \texttt{05\_registro\_alteracoes/registro\_alteracoes.md}.

\begin{table}[htbp]
\centering
\caption{Principais diferenças entre a base do trabalho anterior e a base atualizada}\label{tab:mudancas}
\scriptsize
\begin{tabular}{p{2.6cm}p{5.7cm}p{6.2cm}}
\toprule
Aspecto & Trabalho anterior (P1) & Esta versão (P2)\\
\midrule
Unidades & Espanha; Maceió/AL e Duque de Caxias/RJ & Espanha; Manaus/AM e Parintins/AM (redistribuição do professor)\\
Período & 2013--2022 & 2016--2025 (dez exercícios mais recentes com DCA completa)\\
Fonte municipal & Planilhas do Finbra baixadas manualmente (AL e RJ) & API do Siconfi (DCA, Anexos I-C, I-D e I-E), com coleta registrada\\
Plano de contas & Não tratado; contas lidas por código & Harmonização das três classificações (2016--17, 2018--21, 2022--25) pelo nome da conta, com regra anti-dupla contagem\\
Intraorçamentárias & Incluídas na despesa & Excluídas de receita e despesa\\
FUNDEB & Receita bruta, sem deduções & Receitas líquidas das deduções; FUNDEB recebido entre as transferências\\
Deflator & Fatores pontuais da Calculadora do Cidadão, base 2022, com fator de 1,0054 no próprio ano-base & Média anual do IPCA (Sidra 1737), base 2025 = 1; IPCH anual para a Espanha\\
População & Estimativas até 2021 e Censo 2022 na mesma série; números armazenados como texto & Série harmonizada aos Censos de 2010 e 2022; série oficial mantida para sensibilidade\\
País & FMI (WEO abr./2024), governo geral; ``equilíbrio estrutural'' rotulado de forma imprecisa & Eurostat por subsetor (central, regional, local, seguridade), saldo B.9 e dívida de Maastricht, até 2025\\
Dados ausentes & ``--'' misturado a números; IDEB ausente em anos pares & Ausências explícitas (NA), nunca zero; Parintins 2016 sem DCA registrado\\
Análises & Descritivas, Laffer, regressão simples receita-PIB & Regressões com efeito fixo de município para \textit{flypaper} e esforço tributário, com oito testes de sensibilidade e diagnósticos para amostra pequena; estudo bibliométrico\\
Reprodutibilidade & Planilha & Scripts Python e R que reproduzem bases, tabelas e figuras\\
\bottomrule
\end{tabular}
\par\smallskip
\begin{minipage}{0.97\textwidth}\scriptsize Fonte: elaboração própria a partir da conferência da base de Victor (2026).\end{minipage}
\end{table}

Algumas correções alteram a leitura dos números. A mais importante para a comparação temporal é a da população: misturar estimativas antigas com o Censo de 2022 faria a receita per capita de Parintins saltar artificialmente cerca de 20\% entre 2021 e 2022, apenas pela troca da base populacional. A exclusão das despesas intraorçamentárias reduz a despesa de Manaus em até R\$ 0,9 bilhão em um único ano (2020) e evita que contribuições patronais ao regime próprio sejam contadas como gasto final. Para a Espanha, a substituição do FMI pelo Eurostat corrige um problema conceitual: a variável chamada de ``equilíbrio estrutural'' na base anterior não corresponde ao saldo efetivo das administrações públicas, que em 2020 foi de \esBnoveVinte\% do PIB; a nova base usa a capacidade ou necessidade de financiamento (B.9), conceito empregado nas regras europeias.

\subsection{Manaus e Parintins: financiamento e gasto}

A Tabela~\ref{tab:focos} e a Figura~\ref{fig:composicao} apresentam a evolução das finanças dos dois municípios. Em Manaus, a receita total líquida real per capita passou de R\$ \recManIni{} em 2016 para R\$ \recManFim{} em 2025, crescimento de \recManVar\%. A receita tributária real per capita cresceu mais rapidamente (\tribManVar\%, de R\$ \tribManIni{} para R\$ \tribManFim), puxada pelo ISS, que responde em média por \issShareMan\% da arrecadação tributária, enquanto o IPTU responde por \iptuShareMan\%. Ainda assim, as transferências financiaram entre \depManMin\% e \depManMax\% da receita total no período (média de \depManMedia\%). A composição das transferências reflete a economia da capital: a cota-parte do ICMS, ligada ao valor adicionado da indústria incentivada, responde em média por \icmsShareMan\% das transferências líquidas, o FUNDEB por \fundebShareMan\% e o FPM por apenas \fpmShareMan\%.

\begin{figure}[htbp]
\centering
\caption{Receita real per capita por origem, Manaus e Parintins, 2016--2025 (R\$ de 2025)}\label{fig:composicao}
\includegraphics[width=0.95\textwidth]{fig1_composicao_receita_focos.png}
\par\smallskip
\begin{minipage}{0.95\textwidth}\scriptsize
Fonte: elaboração própria com dados do Siconfi/DCA e do IBGE. Receita corrente própria = receita corrente bruta menos transferências correntes brutas; transferências líquidas incluem as de capital e excluem as deduções para o FUNDEB. Parintins não tem DCA de 2016 no Siconfi.
\end{minipage}
\end{figure}

\begin{table}[htbp]
\centering
\caption{Manaus e Parintins: receitas, despesas e indicadores fiscais, 2016--2025}
\label{tab:focos}
\footnotesize
\begin{tabular}{lrrrrrrr}
\toprule
Ano & \makecell{Receita\\total líq.\\(R\$ pc)} & \makecell{Receita\\tributária\\(R\$ pc)} & \makecell{Transf.\\líquidas\\(R\$ pc)} & \makecell{Despesa\\empenhada\\(R\$ pc)} & \makecell{Resultado\\orçam.\\(R\$ mi)} & \makecell{Depen-\\dência\\(\%)} & \makecell{Rec. trib./\\RCL\\(\%)}\\
\midrule
\multicolumn{8}{l}{\textit{Manaus}}\\
2016 & 3.524,47 & 747,60 & 1.855,85 & 3.382,64 & 273,3 & 52,7 & 24,4\\
2017 & 3.180,48 & 741,13 & 1.999,21 & 3.252,93 & -141,2 & 62,9 & 23,5\\
2018 & 3.522,07 & 887,38 & 2.123,73 & 3.391,50 & 257,3 & 60,3 & 26,4\\
2019 & 3.944,53 & 929,00 & 2.223,45 & 3.826,82 & 234,6 & 56,4 & 26,9\\
2020 & 4.205,71 & 926,14 & 2.428,79 & 4.231,71 & -52,4 & 57,7 & 25,3\\
2021 & 4.178,39 & 954,71 & 2.574,70 & 3.806,60 & 757,9 & 61,6 & 25,5\\
2022 & 4.236,01 & 1.027,30 & 2.668,12 & 4.342,18 & -219,1 & 63,0 & 25,3\\
2023 & 4.609,38 & 1.089,72 & 2.582,93 & 4.493,47 & 250,9 & 56,0 & 26,1\\
2024 & 4.715,63 & 1.198,48 & 2.704,09 & 4.746,61 & -70,6 & 57,3 & 27,2\\
2025 & 4.757,67 & 1.262,24 & 2.739,52 & 4.647,34 & 254,2 & 57,6 & 28,0\\
\multicolumn{8}{l}{\textit{Parintins}}\\
2016 & \multicolumn{7}{c}{DCA não homologada no Siconfi (dado ausente, não imputado)}\\
2017 & 2.597,80 & 170,05 & 2.324,54 & 2.745,85 & -14,6 & 89,5 & 6,6\\
2018 & 3.100,85 & 154,08 & 2.759,83 & 3.128,78 & -2,7 & 89,0 & 5,3\\
2019 & 3.010,26 & 157,02 & 2.727,34 & 2.994,96 & 1,5 & 90,6 & 5,5\\
2020 & 3.789,42 & 162,95 & 3.503,11 & 3.788,54 & 0,1 & 92,4 & 4,6\\
2021 & 4.095,29 & 233,93 & 3.739,29 & 3.960,13 & 13,1 & 91,3 & 6,3\\
2022 & 4.264,52 & 259,50 & 3.806,43 & 4.460,43 & -18,9 & 89,3 & 6,5\\
2023 & 4.878,49 & 265,63 & 4.416,44 & 4.584,00 & 29,2 & 90,5 & 5,7\\
2024 & 6.058,96 & 306,47 & 5.207,27 & 6.227,16 & -17,1 & 85,9 & 5,6\\
2025 & 7.005,42 & 351,06 & 6.210,64 & 6.439,76 & 57,6 & 88,7 & 6,0\\
\bottomrule
\end{tabular}
\par\smallskip
\begin{minipage}{0.97\textwidth}\scriptsize
Fonte: elaboração própria com dados do Siconfi/DCA (STN) e do IBGE (Sidra). Valores em reais de 2025 (IPCA, média anual) e per capita pela população harmonizada (Seção~\ref{sec:tratamentos}). Receita exceto intraorçamentária, líquida das deduções para o FUNDEB; despesa empenhada exceto intraorçamentária. Dependência = transferências correntes e de capital líquidas / receita total líquida.
\end{minipage}
\end{table}


Parintins tem um padrão de outra natureza. As transferências responderam por \depParMin\% a \depParMax\% da receita total (média de \depParMedia\%), e a receita tributária correspondeu, em média, a apenas \autParMedia\% da receita corrente líquida, contra \autManMedia\% em Manaus. Mesmo esse percentual superestima a autonomia: a receita tributária inclui o imposto de renda retido na fonte dos servidores municipais, que pertence ao município por determinação constitucional, mas depende da própria folha de pagamento e não de uma base econômica local. O ISS representa \issSharePar\% da receita tributária e o IPTU, \iptuSharePar\%. As transferências de Parintins são mais diversificadas e mais vinculadas que as de Manaus: o FUNDEB responde por \fundebSharePar\%, o FPM por \fpmSharePar\%, os repasses do SUS por \susSharePar\% e o ICMS por \icmsSharePar\%. Entre 2017 e 2025, as transferências reais per capita cresceram \transfParVar\%, com forte aceleração em 2024 e 2025, impulsionada por repasses do SUS e por transferências de capital, e a despesa real per capita acompanhou esse movimento (\despParVar\%), passando de R\$ \despParIni{} para R\$ \despParFim.

A composição do gasto reforça a diferença. Em Manaus, a despesa real per capita com pessoal cresceu \pessoalManVar\% no período, menos do que a despesa total, enquanto saúde e educação cresceram \saudeManVar\% e \educManVar\%. Em Parintins, a despesa com pessoal cresceu \pessoalParVar\%, a saúde \saudeParVar\% e os investimentos, que partiam de valores muito baixos (R\$ \investParIni{} per capita em 2017), chegaram a R\$ \investParFim{} em 2025 --- uma expansão que coincide com o aumento das transferências de capital. Quanto ao resultado orçamentário, Manaus registrou superávit em \resPosMan{} dos \nAnosMan{} exercícios, com valores reais entre R\$ \resMinMan{} milhões (\anoResMinMan) e R\$ \resMaxMan{} milhões (\anoResMaxMan); Parintins registrou superávit em \resPosPar{} dos \nAnosPar{} exercícios disponíveis. Em ambos os casos, os resultados oscilam em torno do equilíbrio, como esperado de entes sujeitos à LRF e sem acesso relevante a crédito.

\subsection{Dependência de transferências: Manaus, Parintins e Espanha}

A Figura~\ref{fig:dependencia} compara a dependência de transferências dos dois municípios com a dos subsetores espanhóis. A distância entre Manaus e Parintins é estável ao longo de todo o período: a dependência de Parintins nunca ficou abaixo de \depParMin\%, enquanto a de Manaus oscilou entre \depManMin\% e \depManMax\%. A diferença não decorre de oscilações conjunturais, mas da estrutura econômica: o PIB per capita de Manaus em 2023 (R\$ \pibpcManFim) era cerca de três vezes e meia o de Parintins (R\$ \pibpcParFim), e é essa base que sustenta o ISS e a cota-parte do ICMS da capital. A primeira parte da hipótese H1 é, portanto, confirmada.

\begin{figure}[htbp]
\centering
\caption{Transferências recebidas em proporção da receita total: Manaus, Parintins e subsetores da Espanha, 2016--2025}\label{fig:dependencia}
\includegraphics[width=0.9\textwidth]{fig2_dependencia_transferencias.png}
\par\smallskip
\begin{minipage}{0.9\textwidth}\scriptsize
Fonte: elaboração própria com dados do Siconfi, do IBGE e do Eurostat. Para a Espanha, transferências recebidas = D.7 + D.9 sobre a receita total do subsetor (não consolidada). Parintins não tem DCA de 2016 no Siconfi.
\end{minipage}
\end{figure}

Na Espanha, as Comunidades Autônomas receberam em transferências, em média, \esDepCA\% de sua receita, e as entidades locais, \esDepLoc\% (Tabela~\ref{tab:espanha}). A dependência regional espanhola é, portanto, semelhante à de Manaus, mas muito inferior à de Parintins. Duas ressalvas qualificam a comparação. Primeiro, no SEC 2010 as parcelas cedidas do IRPF e do IVA são registradas como impostos regionais, embora as alíquotas e as bases sejam, em larga medida, definidas pelo Estado; uma classificação análoga à brasileira elevaria a dependência regional espanhola. Segundo, as entidades locais espanholas não financiam saúde e educação, que no Brasil absorvem a maior parte das transferências vinculadas. Ainda assim, a segunda parte de H1 é confirmada para Parintins e, parcialmente, para Manaus: os níveis subnacionais espanhóis têm autonomia tributária muito maior (impostos próprios equivalem a \esAutCA\% da receita regional e \esAutLoc\% da receita local, contra \autManMedia\% e \autParMedia\% da receita corrente líquida em Manaus e Parintins), e as entidades locais dependem menos de transferências do que ambos os municípios.

\begin{table}[htbp]
\centering
\caption{Espanha: indicadores fiscais por subsetor das administrações públicas, 2016--2025}
\label{tab:espanha}
\scriptsize
\begin{tabular}{lrrrrrrrr}
\toprule
 & \makecell{Receita\\\% PIB\\(média)} & \makecell{Despesa\\\% PIB\\(média)} & \makecell{Saldo\\\% PIB\\2016} & \makecell{Saldo\\\% PIB\\2025} & \makecell{Dívida\\\% PIB\\2016} & \makecell{Dívida\\\% PIB\\2025} & \makecell{Transf.\\recebidas/\\receita (\%)} & \makecell{Impostos\\próprios/\\receita (\%)}\\
\midrule
Governo geral & 40,7 & 45,0 & -4,2 & -2,4 & 102,0 & 100,5 & 4,9 & 57,0\\
Administração central & 18,4 & 21,5 & -2,5 & -2,0 & 89,9 & 92,5 & 11,0 & 78,5\\
Comunidades Autônomas & 16,0 & 16,5 & -0,8 & -0,4 & 24,7 & 20,2 & 58,9 & 35,4\\
Entidades locais & 6,5 & 6,1 & 0,6 & 0,3 & 2,9 & 1,2 & 44,3 & 47,4\\
Seguridade social & 14,7 & 15,7 & -1,5 & -0,3 & 1,5 & 8,1 & 15,8 & 0,0\\
\bottomrule
\end{tabular}
\par\smallskip
\begin{minipage}{0.97\textwidth}\scriptsize
Fonte: elaboração própria com dados do Eurostat (\texttt{gov\_10a\_main}, \texttt{gov\_10dd\_edpt1}, \texttt{nama\_10\_gdp}). Subsetores não consolidados entre si; transferências recebidas = D.7 + D.9 (inclui transferências intragovernamentais D.73 e D.92); impostos próprios = D.2 + D.5 + D.91, segundo o SEC 2010, que registra como receita regional a parcela cedida do IRPF e do IVA. Saldo = capacidade (+) ou necessidade ($-$) de financiamento (B.9).
\end{minipage}
\end{table}


O contraste mais nítido está na disciplina fiscal (Figura~\ref{fig:espanha}). O déficit das administrações públicas espanholas caiu de \esBnoveGerIni\% do PIB em 2016 para \esBnoveGerFim\% em 2025, depois de atingir \esBnoveVinte\% na pandemia, e a dívida bruta oscilou em torno de 100\% do PIB (\esGDGerIni\% em 2016, \esGDvinte\% em 2020 e \esGDGerFim\% em 2025), muito acima da referência de 60\%. O desequilíbrio, porém, concentra-se na administração central e na seguridade social. As Comunidades Autônomas mantiveram déficits pequenos (média de \esBnoveCAMedia\% do PIB) e reduziram sua dívida de \esGDCAIni\% para \esGDCAFim\% do PIB, e as entidades locais registraram superávit em \esSupLocAnos{} dos dez anos, com dívida que caiu de \esGDLocIni\% para \esGDLocFim\% do PIB. Esse resultado é coerente com H4 e com a evidência de que a autonomia tributária e as regras de estabilidade endurecem a restrição orçamentária subnacional \cite{arespa2022,asatryan2015}. Ele também se aproxima do que se observa em Manaus e Parintins, que, sob a LRF, oscilam em torno do equilíbrio. A diferença é que, na Espanha, a disciplina convive com autonomia tributária, ao passo que em Parintins ela convive com dependência quase total de transferências --- a combinação que \citeonline{rodden2002} associa a pressões por resgate quando as regras de endividamento se afrouxam.

\begin{figure}[htbp]
\centering
\caption{Espanha: saldo e dívida bruta por subsetor das administrações públicas, 2016--2025 (\% do PIB)}\label{fig:espanha}
\includegraphics[width=0.95\textwidth]{fig4_espanha_subsetores.png}
\par\smallskip
\begin{minipage}{0.95\textwidth}\scriptsize
Fonte: elaboração própria com dados do Eurostat (\texttt{gov\_10dd\_edpt1}). Dívida de Maastricht não consolidada entre subsetores.
\end{minipage}
\end{figure}

\subsection{Transferências, renda e gasto em Manaus e Parintins}

A Figura~\ref{fig:dispersao} mostra a relação entre transferências e despesa per capita nos dois municípios, ano a ano. Em Parintins, os pontos sobem praticamente sobre a reta de 45 graus: cada aumento das transferências se converteu em aumento quase equivalente da despesa. Em Manaus, os pontos ficam acima da reta e se espalham mais, porque a despesa também é financiada por receitas próprias, que cresceram no período. A comparação das variações acumuladas confirma essa leitura: entre \propDeMan{} e 2023, a despesa real per capita de Manaus aumentou R\$ \dDespMan{}, para um aumento de R\$ \dTransfMan{} nas transferências (R\$ \propTrMan{} de despesa por real transferido) e de R\$ \dPibMan{} no PIB per capita; em Parintins, entre \propDePar{} e 2023, a despesa aumentou R\$ \dDespPar{} para um aumento de R\$ \dTransfPar{} nas transferências (R\$ \propTrPar{} por real transferido) e de apenas R\$ \dPibPar{} no PIB per capita.

\begin{figure}[htbp]
\centering
\caption{Transferências líquidas e despesa empenhada per capita, Manaus e Parintins, 2016--2025 (mil R\$ de 2025)}\label{fig:dispersao}
\includegraphics[width=0.78\textwidth]{fig3_dispersao_transf_despesa.png}
\par\smallskip
\begin{minipage}{0.78\textwidth}\scriptsize
Fonte: elaboração própria com dados do Siconfi e do IBGE. Cada ponto é um ano (rótulo = dois últimos dígitos); as retas são ajustes lineares por município.
\end{minipage}
\end{figure}

\begin{table}[htbp]
\centering
\caption{Manaus e Parintins: transferências, renda local e gasto per capita, 2016--2023}
\label{tab:reg}
\scriptsize
\setlength{\tabcolsep}{2.6pt}
\begin{tabular}{lccccccccc}
\toprule
 & (1) & (2) & (3) & (4) & (5) & (6) & (7) & (8) & (9)\\
 & Níveis & Tendência & FPM & Log & Dif. & Pop. of. & Manaus & Parintins & Trib.\\
 & Desp. & Desp. & Desp. & ln Desp. & $\Delta$Desp. & Desp. & Desp. & Desp. & Trib.\\
\midrule
Transferências pc ($\beta_1$) & $0,899^{***}$ & $0,550$ &  & $0,839^{***}$ & $0,515$ & $1,040^{***}$ & $0,051$ & $1,054^{**}$ & $0,034$\\
 & (0,119) & (0,418) &  & (0,172) & (0,557) & (0,275) & (1,021) & (0,315) & (0,025)\\[1pt]
FPM pc &  &  & $1,064$ &  &  &  &  &  & \\
 &  &  & (1,040) &  &  &  &  &  & \\[1pt]
Demais transferências pc &  &  & $0,847^{***}$ &  &  &  &  &  & \\
 &  &  & (0,138) &  &  &  &  &  & \\[1pt]
PIB pc ($\beta_2$) & $0,051$ & $-0,024$ & $0,052$ & $0,042$ & $-0,038$ & $0,034$ & $0,129$ & $-0,108$ & $0,031^{***}$\\
 & (0,039) & (0,113) & (0,042) & (0,745) & (0,129) & (0,024) & (0,096) & (0,186) & (0,009)\\[1pt]
Parintins (efeito fixo) & $1,077$ & $-1,833$ & $1,125$ & $-0,300$ &  & $0,303$ &  &  & $0,556$\\
 & (1,661) & (4,463) & (1,846) & (0,955) &  & (1,052) &  &  & (0,386)\\[1pt]
Tendência linear &  & $0,145$ &  &  &  &  &  &  & \\
 &  & (0,178) &  &  &  &  &  &  & \\[1pt]
\midrule
$H_0$: $\beta_1=\beta_2$ ($t$) & 5,84$^{***}$ & 1,75 & 0,96 & 0,89 & 1,04 & 3,41$^{***}$ & -0,07 & 2,35$^{*}$ & --\\
Observações & 15 & 15 & 15 & 15 & 13 & 13 & 8 & 7 & 15\\
$R^2$ ajustado & 0,870 & 0,877 & 0,858 & 0,878 & 0,011 & 0,909 & 0,739 & 0,934 & 0,985\\
\bottomrule
\end{tabular}
\par\smallskip
\begin{minipage}{0.98\textwidth}\scriptsize
Fonte: elaboração própria (código \texttt{03\_analise.py}). MQO com dados empilhados de Manaus e Parintins (município-ano). Variáveis em milhares de reais de 2025 per capita, exceto (4), em logaritmos, e (5), em primeiras diferenças. A coluna (6) usa a população oficial do IBGE (sem 2023); as colunas (7) e (8) estimam cada município separadamente; a coluna (9) tem a receita tributária per capita como variável dependente. Na coluna (3), o teste compara o coeficiente do FPM com o do PIB. Erros-padrão robustos HC3 entre parênteses, com distribuição $t$ de Student. A amostra termina em 2023 porque o PIB municipal ainda não foi divulgado para 2024--2025, e Parintins não tem DCA de 2016. *** $p<0{,}01$; ** $p<0{,}05$; * $p<0{,}10$.
\end{minipage}
\end{table}


A Tabela~\ref{tab:reg} apresenta as estimações. Na especificação principal (coluna 1), com efeito fixo de município, um real adicional de transferências líquidas per capita associa-se a R\$ \bTrUm{} de despesa per capita, enquanto um real adicional de PIB per capita associa-se a R\$ \bPibUm{}, coeficiente não significativo. A hipótese de que as duas propensões são iguais é rejeitada ($t=\wtUm$, $p\,\wpUm$), resultado compatível com o efeito \textit{flypaper}. A ordem de grandeza está em linha com a literatura: a propensão a gastar a renda local fica na faixa de poucos centavos indicada por \citeonline{hines1995}, e a das transferências se aproxima dos valores encontrados para municípios brasileiros \cite{cossio2001,ferreira2019,ferreira2022} e em desenhos quase-experimentais em outros países \cite{dahlberg2008,leduc2017}. O resultado se mantém quando os valores per capita são calculados com a população oficial do IBGE (coluna 6: \bTrSeis{}, com $t=\wtSeis$ para a diferença) e quando se exclui o PIB para usar todos os anos até 2025 (\bTrOito{}, com \nOito{} observações).

A robustez, contudo, é limitada, como se espera com tão poucas observações. Com uma tendência linear comum (coluna 2), o coeficiente das transferências cai para \bTrDois{} e deixa de ser significativo, e a diferença em relação ao PIB passa a ter $p=\wpDois$; em primeiras diferenças (coluna 5), o coeficiente é de \bTrCinco{}, também sem significância. Isso indica que parte da associação em níveis decorre do crescimento simultâneo de transferências e despesas ao longo do tempo, e não apenas da resposta ano a ano. A estimação separada por município (colunas 7 e 8) mostra de onde vem o resultado: em Parintins, o coeficiente das transferências é de \bTrDez{} ($p=\pTrDez$), muito próximo de um real gasto por real recebido; em Manaus, é praticamente nulo (\bTrNove), porque a despesa da capital acompanha também a receita própria. Na especificação logarítmica (coluna 4), a elasticidade do gasto em relação às transferências é de \bTrQuatro{}, mas a do PIB é estimada com tanta imprecisão que a diferença entre as duas não é significativa. A coluna 3 separa o FPM das demais transferências: o coeficiente do FPM (\bFpmTres) é impreciso, porque o FPM per capita varia pouco ao longo do tempo dentro de cada município, enquanto o das demais transferências (\bOutTres) é significativo. A hipótese H2 é, portanto, apenas parcialmente confirmada: o padrão de \textit{flypaper} é claro em Parintins e na especificação em níveis, mas não resiste ao controle de tendência nem às primeiras diferenças.

A coluna 9 estima a Equação~\eqref{eq:trib}. O coeficiente das transferências sobre a receita tributária per capita é pequeno e não significativo (\bTrSete, $p=\pTrSete$), ao passo que o do PIB é positivo e significativo (\bPibSete, $p=\pPibSete$): a arrecadação própria acompanha a economia local, e não as transferências. Não há, portanto, evidência de que transferências maiores se associem a menor arrecadação própria, e a hipótese H3 não é confirmada. O resultado converge com a ausência de deslocamento da receita própria encontrada por \citeonline{litschig2013} e diverge da redução de eficiência arrecadatória documentada por \citeonline{mattos2011} e da substituição observada em municípios com \textit{royalties} \cite{postali2015}. Uma interpretação plausível é que, em um município como Parintins, com receita tributária próxima de \autParMedia\% da receita corrente líquida, não há muito esforço a ser substituído: a margem de ajuste está no gasto, e não na arrecadação.

Os diagnósticos do modelo da coluna 1 ajudam a qualificar esses resultados. Não há indício de heterocedasticidade (Breusch-Pagan com $p=\bpP$), de não normalidade dos resíduos (Shapiro-Wilk com $p=\swP$), de erro de forma funcional (RESET com $p=\resetP$) nem de autocorrelação relevante dos resíduos dentro de cada município (coeficiente AR(1) de \arRho, com $p=\arP$; Durbin-Watson de \dwMan{} em Manaus e \dwPar{} em Parintins). O VIF das transferências é baixo (\vifTr), mas os do PIB (\vifPib) e da variável de Parintins (\vifPar) são muito altos, porque o PIB per capita é quase inteiramente explicado pela diferença entre os dois municípios; é por isso que o coeficiente do PIB é impreciso. A observação mais influente é \cookMun{} em \cookAno{} (distância de Cook de \cookMax{}), ano de forte superávit da capital.

Para a Espanha, a regressão descritiva das despesas regionais e locais per capita sobre as transferências recebidas e o PIB per capita produz coeficientes de \esbTr{} e \esbPib{}, respectivamente ($n=\esN$), com diferença estatisticamente significativa. Com apenas vinte observações e transferências que, nas Comunidades Autônomas, financiam serviços de saúde e educação cujo gasto é em larga medida determinado pelo Estado, o resultado deve ser lido apenas como indício de que a mesma regularidade aparece em um sistema com regras fiscais mais estritas.

\subsection{Discussão: \textit{flypaper}, ilusão fiscal e os dois casos amazônicos}

Os resultados permitem responder à pergunta de pesquisa em três planos. No plano da dependência, Parintins financia cerca de nove de cada dez reais de sua receita com recursos de outras esferas, e mesmo Manaus, com a base própria mais ampla do estado, depende de transferências para mais da metade de sua receita. No plano do comportamento fiscal, o gasto de Parintins acompanha as transferências quase real a real e não responde à renda local, padrão compatível com o efeito \textit{flypaper}; em Manaus, a despesa responde também à receita própria, e o padrão é menos nítido. No plano comparativo, os níveis subnacionais espanhóis têm muito mais autonomia tributária e mantêm resultados próximos do equilíbrio, ao passo que Parintins mantém equilíbrio formal com autonomia mínima.

A ausência de associação entre transferências e esforço tributário, combinada com a forte associação entre transferências e gasto em Parintins, ajuda a discriminar entre as explicações da literatura. Ela é pouco compatível com a versão da ilusão fiscal em que os governos substituem tributos por transferências \cite{mattos2011}, mas é compatível com a versão em que a transferência reduz o preço percebido do bem público e expande o gasto \cite{logan1986,turnbull1998}, com o poder de agenda da burocracia \cite{filimon1982} e, sobretudo, com o argumento do custo da arrecadação: onde explorar a base local é caro, como em municípios insulares de baixa formalização, a resposta às transferências é maior \cite{aragon2013,dahlby2011}. As vinculações constitucionais (saúde, educação e FUNDEB) também aproximam mecanicamente o gasto das transferências \cite{cruz2020}, o que significa que parte da associação medida reflete regras de destinação, e não apenas escolhas locais.

Para a gestão pública, o caso de Manaus mostra que a dependência pode diminuir quando a base econômica sustenta tributos próprios --- a autonomia tributária subiu de \autManIni\% para \autManFim\% da receita corrente líquida no período ---, enquanto o caso de Parintins mostra o limite desse caminho em economias pequenas. A transição para o IBS, que substituirá o ISS e o ICMS, afetará justamente as duas fontes que distinguem Manaus de Parintins\footnote{BRASIL. Emenda Constitucional n.º 132, de 20 de dezembro de 2023 (Reforma Tributária). Disponível em: \url{https://www.planalto.gov.br/ccivil_03/constituicao/emendas/emc/emc132.htm}.}. A experiência espanhola indica que autonomia tributária e regras fiscais podem coexistir com descentralização ampla \cite{arespa2022,asatryan2015}, mas também que a transferência de impostos não elimina a necessidade de mecanismos de nivelamento \cite{zabalza2011}.

\subsection{Dados ausentes, sensibilidade e limitações dos resultados}

Os resultados devem ser lidos à luz de quatro limitações. A primeira é o tamanho da amostra: com 15 observações nas regressões com PIB, os testes têm baixo poder, e a conclusão sobre o efeito \textit{flypaper} depende da especificação, como mostram as colunas 2 e 5 da Tabela~\ref{tab:reg}. A segunda é a ausência da DCA de Parintins em 2016, que não foi imputada; as variações da série do município são calculadas a partir de 2017. A terceira é a cobertura do PIB municipal, que termina em 2023; a especificação sem PIB, que usa os anos até 2025, produz coeficiente semelhante ao da especificação principal. A quarta é a natureza associativa das estimativas: a endogeneidade das transferências e a presença de transferências vinculadas impedem interpretar $\beta_1$ como efeito causal \cite{knight2002}. Mudanças de conceito também afetam a leitura: as três classificações da receita foram harmonizadas pelo nome das contas, mas reclassificações pontuais (por exemplo, o tratamento do IRRF e das complementações da União ao FUNDEB) podem produzir pequenas descontinuidades, registradas no arquivo de auditoria.
